\subsection{Solving a Power Equation} 
If we have an equation in which the only \makeblank[1.5in] appears inside a power function, this is a \makeblank. We solve this by taking an inverse. 
\begin{itemize}
\item If $f(x)=x^p$, then $f^{-1}(x)=x^{1/p}$.
\item If $p$ is rational, write $p=\frac{m}{n}$. Then $f(x)=x^{\makeblanksmall}$ and $f^{-1}(x)=x^{\makeblanksmall}$.
\item If $n$ is even, we have an \makeblank in our equation, so we need to check for \makeblank[2.5in]
\item If $m$ is even, we have an \makeblank in our original equation, so we have \makeblank[2.5in]
\item When we work with power functions, we restrict ourselves to \makeblank[1in] numbers, so we can't take even roots of negative numbers.
\item In many applications, we deal with only positive numbers.
\end{itemize}
\begin{Example}
Solve $3(x-1)^{2/3}=12$
\begin{lpic}{}{9}
\end{lpic}
\end{Example}
